%---
% slug: axosim
% title: AxoSim
% category: product
% eyebrow: NEURAL SIMULATION
% summary: AxoSim replaces detailed neuron simulations with learned models connected through explicit synapses, and one RTX 5090 simulates 65,536 neurons with 4,000 contacts each at 3.38× real-time speed. Learned surrogates thus make detailed neurons practical at population scale.
% author: Davide Wiest
% draft: false
% listed: false
% date: 2026-10-02
% rank: 0.5
% cover_image: /work/covers/axosim.svg
%---
% TODO(v2 retraining): update the abstract's F1 and throughput numbers and the
% figure to the retrained report (copy axo-website public/figures/central-comparison.*).
\documentclass[11pt]{article}
\usepackage{publication-web}
\begin{document}

\publicationtitle{AxoSim: Learned Neuron Models for Population Simulation at Scale}
\publicationauthors{Davide Wiest, Jonathan Sch\"afer}
\publicationaffiliations{Axym Labs}

\begin{publicationhead}
\makepublicationtitle
\end{publicationhead}

\publicationfigure[\linewidth]
  {figures/axosim-central-comparison.png}{/work/axosim-central-comparison.svg}
  {AxoSim-GRU Medium and Mamba Medium in color, alongside Branch-ELM and TCN models and the CoreNEURON reference in gray, compared on inference and history throughput, dataset-relative voltage and subthreshold fidelity, and spike Mean F1.}
  {AxoSim--GRU Medium and Mamba Medium compared with three released Branch--ELM sizes and the smallest and largest released TCN; larger radii indicate better performance. The axes show inference throughput, voltage fidelity (voltage SERA skill as a fraction of the skill of the teacher voltage clipped below action potentials), subthreshold fidelity (RMSE skill outside action-potential windows), spike Mean F1 over 0--5\,ms matching tolerances, and history throughput. The fidelity axes are dataset-relative: each model is scored on held-out traces of its own training data, AxoBench for AxoSim and NeuronIO for Branch--ELM and TCN, with firing rates matched between the datasets. CoreNEURON is the ground-truth detailed simulator, and Branch--ELM is the learned-surrogate predecessor of AxoSim. Raw measurements, fixed mappings and runtime contracts appear in the appendix of the technical report.}

\begin{publicationsummary}
\publicationsummarypart{Problem}{Detailed neuron simulators model each neuron's branches and ion channels, but are costly at population scale. Learned models reduce this cost; in a connected network, however, each predicted spike becomes another neuron's input, so errors change the model's future inputs.}
\publicationsummarypart{Approach}{AxoSim connects learned neuron models, which account for each cell's shape, through synapses with delays and adjustable strengths, and keeps the population trainable with gradients.}
\publicationsummarypart{Evidence}{Its GRU neuron model reproduces the simulator's spikes with a mean F1 of 0.61, against up to 0.45 for released Branch-ELM and TCN models on their own training data, at 16.8 million neuron-steps per second, over 20 times Branch-ELM's throughput. One RTX~5090 simulates 65,536 neurons with 4,000 synaptic contacts each at 3.38$\times$ real-time speed and trains 20,000 neurons with 2,000 contacts each in real time.}
\publicationsummarypart{Adaptation}{When the timing of spoken digits shifts, adapting morphology and synapse parameters raises accuracy by $6.46\pm1.20$ percentage points across six independently trained models.}
\publicationsummarypart{Data}{A released set of 50,000 validated population traces from 100 recurrent networks supports evaluation in connected circuits.}
\publicationsummarypart{Significance}{These experiments thus separate how well a model imitates a single neuron from how it behaves and adapts in a recurrent network.}
\end{publicationsummary}

\publicationmodels{https://huggingface.co/Axym-Labs/AxoSim}
\publicationdata{https://huggingface.co/datasets/Axym-Labs/axobench-population}
\publicationcode{https://github.com/Axym-Labs/axosim}
\publicationdocs{https://axo.axym.org}
\publicationpackage{https://pypi.org/project/axosim/}
\publicationdemo{/work/axosim-connectome-demo/}
\publicationcorrespondence{dw@axym.org}
\makepublicationlinks

\begin{publicationcitation}
\begin{verbatim}
@techreport{wiest2026axosim,
  title       = {AxoSim: Learned Neuron Models for
                 Population Simulation at Scale},
  author      = {Wiest, Davide and Sch{\"a}fer, Jonathan},
  institution = {Axym Labs},
  year        = {2026},
  url         = {https://axym.org/work/axosim/}
}
\end{verbatim}
\end{publicationcitation}

\end{document}
